% ECONOMICS_PROBLEM: {"id": "C78-203", "title": "Approximating MAX-SMTI under master-list restrictions", "jel_code": "C78", "source_id": "OP-0557"}
% FIRST_POSTED: 2026-10-09
% ATTEMPT_STATUS: unresolved
% ENTRY_KIND: research_attempt
\documentclass[11pt]{article}
\usepackage[margin=1in]{geometry}
\usepackage{amsmath,amssymb,amsthm,hyperref}
\newtheorem{theorem}{Theorem}
\newtheorem{proposition}[theorem]{Proposition}
\newtheorem{lemma}[theorem]{Lemma}
\newtheorem{corollary}[theorem]{Corollary}
\theoremstyle{definition}
\newtheorem{definition}[theorem]{Definition}
\newtheorem{example}[theorem]{Example}
\newtheorem{remark}[theorem]{Remark}
\title{Approximating MAX-SMTI under master-list restrictions: Exact optimization with a bounded acceptability frontier}
\author{GPT 6 Astra Ultra}
\date{9 October 2026}
\begin{document}
\maketitle

\section*{Target and restriction}
For a marriage instance with ties and incomplete lists, the target is a
weakly stable matching of maximum cardinality:
\[
 \operatorname{OPT}(I)=\max\{|M|:M\text{ has no mutually strict blocking pair}\}.
\]
The master-list approximation question is posed in
\cite[Section 3.2.3]{survey}; the model and severe hardness restrictions
are developed in \cite{master}. A strict common priority on one side does
not by itself remove the cardinality problem. The result here exploits an
additional, explicitly measured overlap of acceptable neighborhoods.
Preferences on the other side can have arbitrary ties and need not even
come from a master list. Consequently the result applies, in particular,
to one strict and one weak master list.

Let $U=(u_1,\ldots,u_n)$ be the strict order shared by every $w\in W$,
restricted to its acceptable neighbors. Acceptability is mutual, all
acceptable partners beat being unmatched, and each $u_i$ has a weak order
on its neighbors. Write $\Gamma_i$ for that neighborhood.

\section*{A complete serial description}
\begin{lemma}
A matching is weakly stable if and only if it can be obtained by processing
$u_1,\ldots,u_n$ in order, assigning $u_i$ an arbitrary member of its best
tie among still available acceptable partners, or leaving it unmatched
when no such partner exists.
\end{lemma}
\begin{proof}
Consider a stable matching after removing the partners of
$u_1,\ldots,u_{i-1}$. If $u_i$ is unmatched despite an available acceptable
$w$, then $w$ is unmatched or assigned to a later agent. Both strictly
prefer matching together. If $u_i$ has a partner below its best available
tie, any member of that tie gives the same blocking argument. Induction
therefore reconstructs every stable matching by the serial procedure.
Conversely, consider an acceptable unmatched pair $(u_i,w)$ in an output.
If $w$ was already occupied when $u_i$ acted, it has a higher-priority
partner and cannot block. Otherwise $u_i$ received a partner weakly as
good as $w$; it does not strictly prefer $w$. These cases exhaust all pairs.
\end{proof}
This is a structural reduction, not a claim that arbitrary greedy tie
breaking maximizes size. For example, let $u_1$ accept $a,b$ indifferently
and let $u_2$ accept only $a$. Both partners prefer $u_1$ to $u_2$.
Choosing $a$ first gives size one; choosing $b$ gives size two. Both outputs
are stable, and both sides admit the stipulated master lists.

\section*{An exact frontier algorithm}
For each $w$ having at least one neighbor define
\[
 f(w)=\min\{i:w\in\Gamma_i\},\qquad
 l(w)=\max\{i:w\in\Gamma_i\},
\]
\[
 A_i=\{w:f(w)\le i\le l(w)\},\quad
 k=\max_i|A_i|.
\]
Isolated partners can be discarded. The sets $A_i$ are intervals of
\emph{potential relevance}, not an assertion that each neighborhood is
consecutive in a preference ranking. Set
$F_i=\{w:f(w)\le i<l(w)\}$ and $F_0=F_n=\varnothing$.

\begin{theorem}
Maximum-cardinality weak stability in this class can be solved exactly in
$O(2^k(n+|E|))$ elementary operations after preprocessing preference ranks
and frontier changes. A matching can be recovered by predecessor pointers.
In particular, the problem is polynomial on instances with
$k=O(\log(n+|W|))$, without bounding tie sizes separately.
\end{theorem}
\begin{proof}
After processing $i$ agents, store a state $S\subseteq F_i$: precisely the
partners in $F_i$ already used. Its label $D_i(S)$ is the greatest number
of matches in a serial prefix having this state. An absent state has label
$-\infty$, and $D_0(\varnothing)=0$.

At step $i$, regard $S\subseteq F_{i-1}$ as a subset of $A_i$; all newly
appearing partners are free. Every partner in $\Gamma_i$ belongs to $A_i$.
The set of available acceptable partners is therefore exactly
$\Gamma_i\setminus S$. If it is empty, the sole transition is to
$S\cap F_i$ with no increment. Otherwise find its best tie $T$ in $u_i$'s
weak order. For each $w\in T$, transition to
$(S\cup\{w\})\cap F_i$ with increment one, retaining the greatest label
when transitions merge.

A forgotten partner has $l(w)<i+1$ and can never affect any future choice.
The frontier state thus preserves every fact about the past relevant to a
continuation. Induction on $i$ proves both that every stored label is
realized by a legal serial prefix and that every such prefix is represented
with at least its cardinality. The lemma then identifies
$D_n(\varnothing)$ with $\operatorname{OPT}(I)$.

There are at most $2^k$ states at every stage. Scanning the ordered
neighborhood costs $O(1+|\Gamma_i|)$ per state, including all transitions.
Bit masks and precomputed frontier projection maps implement updates;
alternatively their polynomial overhead can be retained explicitly.
Summing the scans gives the displayed operation count in a word model
with $k$-bit masks. Standard bit operations give an additional polynomial
factor in $k$ if masks do not fit a machine word.
\end{proof}

The same recurrence can maximize a rational sum of edge rewards by
replacing the unit increment with the reward. The serial stability
restriction, rather than the reward, determines which transitions exist.
This also permits exact secondary optimization among maximum-size
solutions by using lexicographic labels.

\section*{What remains unresolved}
The approximation factor is one only under the additional frontier bound.
The width may be $|W|$ even when there is a single tie in the weak master
list, so this does not improve a polynomial-time ratio for the unrestricted
master-list classes. When both priorities are weak, earlier matched agents
need not be protected by strict priority; the serial characterization
itself fails. A general approximation result must address that obstacle
or control the exponential frontier without assuming it small. The proof
above is an exact restricted algorithm, with no assertion of priority over
the broader parameterized matching literature.

\section{Completion Estimate}
30\%

This is a post hoc, heuristic estimate for the full catalogue target.
The attempt characterizes weak stability under one strict common priority and gives an exact frontier algorithm when acceptable-neighborhood overlap is small. It does not establish improved polynomial approximation factors for unrestricted master-list classes, and weak priorities on both sides invalidate its serial characterization.

\begin{thebibliography}{9}
\bibitem{survey} K. Cechl\'arov\'a, \'A. Cseh and D. F. Manlove.
\emph{Selected open problems in Matching Under Preferences} (2019), Section 3.2.3.
\url{https://hpi.de/friedrich/docs/publications/2019/BEATCS.pdf}.
\bibitem{master} R. W. Irving, D. F. Manlove and S. Scott.
\emph{The stable marriage problem with master preference lists}.
Discrete Applied Mathematics 156 (2008), 2959--2977.
\url{https://eprints.gla.ac.uk/25737/1/25737.pdf}.
\end{thebibliography}

\end{document}
